\section{Two points and direction vectors}

Many line problems begin with two points. This is natural: in Euclidean geometry, two distinct points determine exactly one line. In analytic geometry we translate this statement into coordinates and vectors.

The translation is simple. Two points give a displacement. That displacement gives a direction.

\subsection*{Building the direction vector}

Let
\[
A=(a_1,a_2),
\qquad
B=(b_1,b_2),
\]
and suppose \(A\ne B\). The vector from \(A\) to \(B\) is
\[
\overrightarrow{AB}=B-A=(b_1-a_1,\,b_2-a_2).
\]
This vector is non-zero because the points are distinct. It is therefore a valid direction vector for the line through \(A\) and \(B\).

The line can then be written as
\[
P=A+t(B-A),
\qquad t\in\mathbb{R}.
\]

\begin{examplebox}
Find a parametric equation of the line through
\[
A=(1,2),
\qquad
B=(4,-1).
\]
First compute the direction vector:
\[
B-A=(4-1,\,-1-2)=(3,-3).
\]
Therefore the line is
\[
P=(1,2)+t(3,-3).
\]
In coordinate form,
\[
x=1+3t,
\qquad
y=2-3t.
\]
\end{examplebox}

This procedure should be read as a chain of meanings. The two points determine a displacement. The displacement gives a direction. A point and a direction determine the line.

\subsection*{The order of the points}

If we reverse the order of the two points, the direction vector changes sign:
\[
A-B=-(B-A).
\]
This does not change the line. It only reverses the orientation of the parametrization.

Using the previous example, we found
\[
B-A=(3,-3).
\]
If we start from \(B\) instead, then
\[
A-B=(-3,3).
\]
The line may also be written as
\[
P=(4,-1)+s(-3,3).
\]
This is the same line, but the parameter moves along it in the opposite orientation.

\begin{remarkbox}
Changing the orientation of the direction vector does not change the geometric line. It changes how the line is traced by the parameter.
\end{remarkbox}

\subsection*{Non-unique direction vectors}

A direction vector for a line is never unique. If \(v\) is a direction vector, then every non-zero scalar multiple \(\lambda v\) is also a direction vector for the same line.

For instance, the vectors
\[
(3,-3),
\qquad
(1,-1),
\qquad
(-2,2)
\]
all describe the same direction. They have different lengths and possibly different orientations, but they are parallel.

\begin{examplebox}
The line
\[
P=(1,2)+t(3,-3)
\]
can also be written as
\[
P=(1,2)+s(1,-1).
\]
The second version uses a shorter direction vector. Both equations describe the same set of points, because \((3,-3)=3(1,-1)\).
\end{examplebox}

This is often useful. We may simplify a direction vector by dividing out a common factor. The line remains the same, but the notation becomes easier to read.

\subsection*{Finding a point from a parametrization}

A parametric line already contains a point. In
\[
P=A+t v,
\]
the point \(A\) appears when \(t=0\). But every value of \(t\) gives another point on the same line. We may choose any of them as a new anchor point.

\begin{examplebox}
Consider
\[
P=(2,-1)+t(4,2).
\]
For \(t=\frac12\), we get
\[
P=(2,-1)+\frac12(4,2)=(4,0).
\]
Thus \((4,0)\) is also a point on the line. We could write the same line as
\[
P=(4,0)+s(4,2).
\]
\end{examplebox}

This flexibility is not a problem. It is part of the geometry. A line has infinitely many points and infinitely many possible non-zero direction vectors.

\subsection*{When the two points are the same}

If the two given points are equal, then they do not determine a unique line. The vector between them is the zero vector:
\[
B-A=0.
\]
The zero vector cannot be a direction vector, because it gives no direction.

This is a common hidden issue in problems with parameters. If two points depend on a parameter, there may be special parameter values for which the points coincide. Those values must be checked separately.

\begin{summarybox}
When a line is given by two points, ask:
\begin{itemize}
\item Are the two points distinct?
\item What is the vector from one point to the other?
\item Can the direction vector be simplified?
\item Does reversing the order only reverse orientation?
\item Are there parameter values for which the two points might coincide?
\end{itemize}
\end{summarybox}

The formula \(P=A+t(B-A)\) is not just a shortcut. It expresses the idea that a line is obtained by moving from one point in the direction of another.